\subsection{Integrality of the original native polynomial product model}
\begin{lemma}[The entire actual coefficient extension ideal]
\label{lem:actual-polynomial-coefficient-ideal-membership}
\lean{EnvelopeProperty.mem_polynomialCoefficientIdeal_iff}
\leanok
\uses{polynomial:coefficient-variable-interchange}
Let $R$ be ANY commutative ring, $\sigma$ ANY variable set, and
$J\subseteq R[\sigma]$ ANY original ideal. Let
$q:R[\sigma]\to R[T][\sigma]$ be the ORIGINAL coefficient map induced
by $R\to R[T]$, and let $e:R[T][\sigma]\simeq R[\sigma][T]$ be the
EXISTING polynomial-coefficient interchange. For EVERY original
$F\in R[T][\sigma]$,
\[F\in q(J)R[T][\sigma]\quad\Longleftrightarrow\quad
  \forall m,\ \operatorname{coeff}_m(eF)\in J.\]
Here the extension is the entire actual nonradical ideal image.
In particular $q(a)\in q(J)R[T][\sigma]$ if and only if $a\in J$.
No field, finite variable set, reducedness, Noetherian or finiteness
assumption is imposed. Original nilpotents are retained.
\end{lemma}
\begin{proof}
The EXISTING actual interchange satisfies $e\circ q=C$, the original
univariate constant inclusion: the existing formulas on the SAME original
coefficient constants and each original multivariable generator prove this
whole ring-map equality by polynomial extensionality. Thus the actual
ideal image under $e$ is EXACTLY the original ideal $J$ extended along $C$,
by functoriality of ideal images. The original $e$ is a RING EQUIVALENCE, so its actual ideal image is
exactly the set image (surjectivity is used); its injectivity then gives
that membership in the SAME ideal is equivalent to membership of $eF$
in that actual ideal image. Equivalently use the existing ideal-equivalence
image/comap membership identity.
Apply the EXISTING Mathlib entire-ideal coefficient criterion for extension
along $C$. For $F=q(a)$ the polynomial $eF=C(a)$ has coefficient $a$ in
degree zero and zero in every other degree; the criterion gives the final
membership reflection. No new interchange or quotient construction is made.
\end{proof}

\begin{lemma}[Actual prime coefficient extension]
\label{lem:actual-polynomial-coefficient-prime-extension}
\lean{EnvelopeProperty.isPrime_polynomialCoefficientIdeal}
\leanok
\uses{lem:actual-polynomial-coefficient-ideal-membership,polynomial:coefficient-variable-interchange}
Under the preceding arbitrary-ring hypotheses, if the ORIGINAL ideal $J$
is prime, then its actual coefficient extension $q(J)R[T][\sigma]$ is prime.
\end{lemma}
\begin{proof}
Reuse Mathlib's existing primality of $J$ extended to $R[\sigma][T]$
along $C$. The actual image of the original extended ideal under $e$
is exactly this prime ideal, by the preceding whole-map calculation.
Its inverse image under the original ring equivalence is prime and is
literally the original coefficient extension. No prime ideal, field,
geometric-integrality or reduced quotient certificate is an input.
\end{proof}

\begin{theorem}[The original native variety times the polynomial line is integral]
\label{thm:actual-native-polynomial-product-integral}
\lean{EnvelopeProperty.ProjSubvariety.nativeProductModel_isIntegral}
\leanok
\uses{lem:actual-polynomial-coefficient-ideal-membership,lem:actual-polynomial-coefficient-prime-extension,lem:original-prime-projective-kernel-integral,def:pl-native-variety-product,projective:actual-polynomial-line-product,projective:actual-product-model-ideal-functions}
Let $k$ be ANY field and $X$ ANY original projective subvariety. Its EXISTING
\texttt{NativeProductModel}, with its original scheme structure, is integral.
No algebraic closure, perfection, characteristic, geometric integrality or
caller-supplied integrality assertion is assumed.
\end{theorem}
\begin{proof}
Let $J$ be the ORIGINAL homogeneous prime of $X$, and choose the ORIGINAL
relevant degree-one coordinate $x_j\notin J$. Take the ACTUAL homogeneous
coefficient extension $J'$ over $k[T]$. Its underlying ideal is prime by
the preceding theorem, and its original image coordinate $q(x_j)$ remains
outside $J'$ by the entire coefficient-ideal criterion. This original
coordinate is still homogeneous of degree one for the ACTUAL original
polynomial projective grading over $k[T]$.
Apply the already constructed general homogeneous-prime quotient theorem
to conclude that the literal kernel closed subscheme for $J'$ inside the
original polynomial projective scheme is integral. This reuses the genuine
graded-domain integral-Proj owner; no generic Proj proof is reproduced.

The EXISTING exact theorem \texttt{nativeProductCutIdeal\_eq} identifies
the ORIGINAL product cutting ideal with the actual pullback of that kernel
ideal along the inverse of the EXISTING ambient polynomial-line product
scheme isomorphism. The genuine geometric ideal-comap closed-subscheme
isomorphism followed by the corresponding pullback projection is an actual
isomorphism onto the same original kernel subscheme. Transport integrality
along that actual isomorphism. Its source is definitionally the EXISTING
\texttt{NativeProductModel}; no change of homogeneous ideal, saturation,
reduced support, variety, polynomial coordinate or structure map occurs.
\end{proof}

The existing native subvariety integrality is a separate reused owner.
The author's existing whole Hironaka--DeJong blowup isomorphism, together
with the actual integrality consumer proved separately, now transfers the
existing DeJong global integrality theorem to the SAME original Hironaka
blowup. Finite normalization, the normalized action and relative nefness
remain separate obligations. This slice changes neither original blowup
factory and does not claim LocalPositivity.
